\documentclass[11pt]{article}

% ----------------------------------------------------------------------
%  The Value-Creation Bridge — client-facing methodology one-pager
%  Draft concept prepared by the ENT 105 team for Creo Advisors.
%  Compile: pdflatex value-creation-bridge.tex
% ----------------------------------------------------------------------

\usepackage[utf8]{inputenc}
\usepackage[T1]{fontenc}
\usepackage{textcomp}
\usepackage[margin=0.9in]{geometry}
\usepackage{booktabs}
\usepackage{array}
\usepackage{enumitem}
\usepackage{titlesec}
\usepackage{xcolor}
\usepackage{tcolorbox}
\usepackage[hidelinks]{hyperref}

\definecolor{creonavy}{RGB}{17,38,73}
\definecolor{creoaccent}{RGB}{27,94,140}

\titleformat{\section}{\large\bfseries\color{creonavy}}{}{0em}{}
\titlespacing*{\section}{0pt}{1.5ex plus .3ex}{0.7ex}
\setlist[itemize]{leftmargin=1.3em,topsep=2pt,itemsep=2.5pt}
\setlength{\parskip}{4pt plus 1pt minus 1pt}
\setlength{\parindent}{0pt}

\newtcolorbox{hook}{colback=creonavy!6,colframe=creonavy,
  boxrule=0.8pt,arc=2pt,left=8pt,right=8pt,top=6pt,bottom=6pt}
\newtcolorbox{sowhat}{colback=creoaccent!8,colframe=creoaccent,
  boxrule=0.6pt,arc=2pt,left=8pt,right=8pt,top=5pt,bottom=5pt}

\newcommand{\created}{\textcolor{creoaccent}{\textbf{Created}}}
\newcommand{\transferred}{\textcolor{gray!55!black}{\textbf{Transferred}}}

\begin{document}

\begin{center}
{\LARGE\bfseries\color{creonavy} The Value-Creation Bridge}\\[4pt]
{\large\color{creoaccent} How much of your return did you \emph{create} --- and how much was
\emph{transferred} to you?}\\[5pt]
{\normalsize Separating operational value from financial engineering --- with no double-counting.}
\end{center}

\vspace{2pt}\hrule\vspace{8pt}

\begin{hook}
\textbf{The question every LP is now asking.} With multiple expansion gone and leverage
expensive, a strong return is no longer self-evidently a \emph{good} return. The question is
how much of it you actually \textbf{built} --- real, repeatable operational value --- versus
how much was \textbf{handed to you} by a richer exit multiple, cheap leverage, or a tax
structure. The Value-Creation Bridge answers that precisely, and then shows exactly where the
value you built came from.
\end{hook}

\section{What it answers}
Two questions, in order:
\begin{enumerate}[leftmargin=1.6em]
  \item \textbf{Created or transferred?} Of the total value realized over the hold, how much
  is genuine operational value the business \emph{created}, and how much is value
  \emph{transferred} into (or out of) your equity --- multiple re-rating, deleveraging, tax
  shields, timing?
  \item \textbf{Where did the created value come from?} The portion you genuinely created,
  attributed cleanly across the operating levers --- \textbf{pricing}, \textbf{volume}, and
  \textbf{cost productivity} --- with capital efficiency tying the operations to cash.
\end{enumerate}

\section{How it works}
\textbf{Split 1 --- Created vs.\ Transferred.} Take the whole system --- the company and every
counterparty (buyers, lenders, the tax authority, the next owner). Value the business
\emph{made} shows up as the total pie growing; value merely \emph{moved} between
parties leaves the total unchanged and only reshuffles who holds it. Those are two
non-overlapping pools, and together they account for the entire gain --- nothing falls
between them.

\smallskip
\textbf{Split 2 --- Attribution within Created.} The created pool is then split across pricing,
volume, and cost productivity so that the pieces \emph{add up exactly} to the total created,
with no leftover ``interaction'' or ``other'' line and no lever double-counted.

\section{Why you can trust the numbers}
This is what separates the Bridge from a conventional value-bridge slide:
\begin{itemize}
  \item \textbf{It sums exactly to the whole.} The pieces reconcile to the realized gain to
  the dollar --- it is grounded in conservation-based accounting (the same balance principle
  as double-entry bookkeeping), so the books balancing \emph{is} the decomposition being
  exact.
  \item \textbf{No double-counting.} Created and transferred do not overlap, and each operating
  lever is credited once. The cross-terms that usually get dumped into a residual are
  allocated by a single, symmetric, defensible rule.
  \item \textbf{Every assumption is on the page.} The math fixes the \emph{form}; the
  judgment calls --- where the system boundary sits, how flows are classified, how price is
  separated from volume --- are stated explicitly, so the answer is transparent and
  contestable rather than a black box.
\end{itemize}

\section{What we need / what you get}
\begin{center}
\renewcommand{\arraystretch}{1.25}\small
\begin{tabular}{p{6.6cm} p{6.6cm}}
\toprule
\textbf{Inputs (from you)} & \textbf{Outputs (from us)} \\
\midrule
Entry and exit financials over the hold; deal structure (multiple, leverage, taxes); basic
price and volume data per major line. &
A clean Created-vs-Transferred bridge; a one-line ``real vs.\ engineered'' verdict; the
created value attributed to pricing / volume / cost; and the implication for the next hold's
value-creation plan. \\
\bottomrule
\end{tabular}
\end{center}

% ======================================================================
\section{Worked example (illustrative)}
% ======================================================================
A sponsor exits a business after five years at a headline \textbf{3.1$\times$ MOIC} ---
equity in at \$400M, out at \$1{,}250M, a \textbf{\$850M gain}. A celebrated outcome. The
Bridge asks: how much did the company actually create?

\smallskip
\textbf{Split 1 --- of the \$850M gain, how much was created vs.\ transferred?}
\begin{center}
\renewcommand{\arraystretch}{1.25}\small
\begin{tabular}{l r l}
\toprule
\textbf{Source of the equity gain} & \textbf{\$M} & \textbf{Bucket} \\
\midrule
Operational EBITDA growth (\$100M $\rightarrow$ \$150M, at entry multiple) & 500 & \created \\
Multiple expansion (10.0$\times$ $\rightarrow$ 11.0$\times$) & 150 & \transferred \\
Deleveraging (net debt \$600M $\rightarrow$ \$400M) & 200 & \transferred \\
\midrule
\textbf{Total equity gain} & \textbf{850} & \\
\bottomrule
\end{tabular}
\end{center}

\begin{sowhat}
\textbf{So what.} Only \textbf{\$500M --- 59\% --- was value the company created.} The other
\textbf{41\% (\$350M) was transferred in}: a richer exit multiple (\$150M, paid by the next
buyer) and deleveraging (\$200M --- value moving from debt to equity as the debt is retired;
enterprise value is unchanged, so it is a capital-structure transfer, not new value). Neither
is operational alpha, and neither is guaranteed to repeat. Value the business at the
\emph{entry} multiple on the \emph{entry} balance sheet and exit equity is about \$900M: the
headline 3.1$\times$ is really a \textbf{2.25$\times$ of created value} wrapped in financial
engineering.
\end{sowhat}

\textbf{Split 2 --- of the \$500M created, where did it come from?}
\begin{center}
\renewcommand{\arraystretch}{1.25}\small
\begin{tabular}{l r r}
\toprule
\textbf{Operating lever} & \textbf{\$M} & \textbf{Share} \\
\midrule
Pricing (net realized price, $+8.0\%$)        &  95 & 19\% \\
Volume ($+18.7\%$)                            & 211 & 42\% \\
Cost productivity (margin, $+17.0\%$)         & 194 & 39\% \\
\midrule
\textbf{Total created} & \textbf{500} & \textbf{100\%} \\
\bottomrule
\end{tabular}
\end{center}

{\footnotesize The three levers compound to the $+50\%$ EBITDA growth
($1.080\times1.187\times1.170=1.50$, i.e.\ \$100M$\rightarrow$\$150M), and the dollars are the
exact log decomposition of that growth --- so they sum to the created \$500M with no
``interaction'' or ``other'' residual.}

\begin{sowhat}
\textbf{So what.} Of the value the company \emph{built}, volume and cost productivity did
\textasciitilde{}80\% of the work; \textbf{pricing contributed just 19\%.} For a business with
real brand and switching costs, that is the loudest signal in the analysis: \textbf{pricing is
the under-exploited lever} --- and the most credible source of value for the next hold.
\end{sowhat}

\textbf{Assumptions --- stated, because that is the point.} Multiple expansion is treated as
transferred (no evidence of structural de-risking behind the re-rating); EBITDA growth is
valued at the \emph{entry} multiple so growth is not conflated with re-rating; deleveraging is
classified as a within-capital-structure transfer (enterprise value is unchanged by the
paydown --- a sponsor who credits the cash flow that funded it to operations would reclassify
it, and the bridge updates); price is separated from volume using the company's own realized
price index. Change any of these and the bridge updates transparently
--- the rigor is exactly as good as the assumptions, and the assumptions are visible.

\vspace{6pt}\hrule\vspace{5pt}
{\footnotesize\itshape Draft methodology concept prepared by the ENT 105 team for Creo
Advisors. Figures are illustrative. Underlying decomposition grounded in conservation-based
accounting; assumptions are client-specific and set at the engagement's scoping stage.}

\end{document}
